\documentclass{article}
\usepackage{float}
% Language setting
\usepackage[german]{babel}
\usepackage[table]{xcolor}
% Set page size and margins
% Replace `letterpaper' with `a4paper' for UK/EU standard size
\usepackage{multicol}
\usepackage[letterpaper,top=2cm,bottom=2cm,left=3cm,right=3cm,marginparwidth=1.75cm]{geometry}
\usepackage[dvipsnames]{xcolor}
% Useful packages
\usepackage{amsmath}
\usepackage{graphicx}
\usepackage[colorlinks=true, allcolors=blue]{hyperref}
\usepackage{subcaption}
\usepackage{graphicx}
\title{Test - Fragen auf Kolloquien} 
\author{Lil'Jana}
\usepackage{enumitem,amssymb}
\newlist{todolist}{itemize}{2}
\setlist[todolist]{label=$\square$}
\usepackage{pifont}




\begin{document}
\maketitle
\subsection*{Die DNA, die im Gentechnik-Labor bearbeitet wird, wird meist zunächst zerkleinert durch ...}
\begin{todolist}
\item Ligasen  
\item  Restriktionsenzyme
\item DNA-Polymerase  
\item starke Säure  
\item Röntgenstrahlen  
\end{todolist}
\subsection*{Was sind die wesentlichen Unterschiede zwischen DNA und RNA?}
\begin{itemize}
    \item 
    \item 
\end{itemize}
\subsection*{Welche Eigenschaften sollte ein Vektor mindestens haben?}
\begin{itemize}
    \item 
    \item 
    \item    
\end{itemize}
\subsection*{Wofür steht die Abkürzung 'PCR'?}
\begin{itemize}
    \item 
\end{itemize}

\subsection*{Das fehlerhafte Protein bei der (erblichen Form der) Hypercholesterinämie ist ...}
\begin{todolist}
\item ein Enzym aus dem Aminosäure-Stoffwechsel.  
\item ein Enzym aus dem Lipidstoffwechsel.   
\item das Hämoglobin. 
\item  der LDL-Rezeptor. 
\item ein Enzym aus dem Kohlenhydrat-Stoffwechsel. 
\end{todolist} 
\subsection*{Wie lauten die vier fehlenden Bezeichnungen?}
\begin{figure}[H] 
    \centering
    \includegraphics[width=10cm]{bilder/nukleotif-fehlend.png}
    \caption{}
    \label{fig:my_label}
\end{figure}
\subsection*{Die Funktion von Restriktionsenzymen in Bakterien ist ....}
\begin{todolist} 
\item die eigene DNA zu schneiden, um Rekombination zu ermöglichen  
\item  fremde DNA zu zerstören
\item fremde Proteine zu zerstören  
\item DNA zu synthetisieren  
\item Plasmide zu zerschneiden  
\end{todolist}
\subsection*{Homologe Chromosomen unterscheiden sich häufig in den Merkmals-Informationen, die sie tragen. Sie ...}
\begin{todolist} % %
\item sind deswegen unterschiedlich lang.  
\item sind deswegen beide in einer Keimzelle enthalten.  
\item sind in diploiden, somatischen Zellen nicht beide vorhanden.  
\item  enthalten deswegen öfter unterschiedliche Allele.
\end{todolist}% %
\subsection*{Die tRNA, die die Aminosäure Valin translatiert, hat das Anticodon CAG.
Wie lautet das dazugehörige mRNA-Triplett?}
\begin{todolist}
\item CAG  
\item GAC  
\item GTC  
\item TUG  
\item  GUC
\end{todolist}
\subsection*{In welchem Meiosestadium lagern sich die homologen Chromosomen zusammen, und findet "crossover" statt?}
\begin{todolist} % %
\item Prophase II  
\item  Prophase I
\item Metaphase II  
\item Metaphase I  
\end{todolist}
\subsection*{Wie viele Typen von Keimzellen kann ein diploider Organismus mit n = 4
unabhängig vom "crossover" bilden?}
Anzahl der Chromosomenpaare:  \newline
Mögliche Kombinationen: 
\begin{todolist} % %
\item 2
\item  4
\item 8
\item 16
\item 32
\item 64
\end{todolist}
\subsection*{Welche der genannten Aussagen ist/sind FALSCH?}
\begin{todolist} % %
\item Haploide Zellen können sich durch Mitose teilen.  
\item  Haploide Zellen können sich durch Meiose teilen.  
\item Diploide Zellen können sich durch Mitose teilen.  
\item Diploide Zellen können sich durch Meiose teilen.  
\end{todolist}
\subsection*{Der Mensch hat 46 Chromosomen. Diese Anzahl (menschlicher)
Chromosomen findet man in ...}
\begin{todolist} % %
\item allen Zellen des Körpers.  
\item allen Eizellen und Spermien.  
\item allen Zellen der Keimbahn nach der Meiose I.  
\item  allen somatischen Zellen.
\end{todolist}
\subsection*{Ein Karyogramm wird hergestellt, in dem man ...}
\begin{todolist}% % 
\item den Phänotyp eines Individuums betrachtet.  
\item den Genotyp eines Individuums analysiert.  
\item eine einzige Chromosomenkombination aus einer Keimzelle
betrachtet.  
\item die Blutgruppe eines Individuums bestimmt.  
\item die homologen Chromosomen einer Zelle nach Anzahl, Größe
und Gestalt ordnet.
\end{todolist}
\subsection*{Durch Alloploidie soll eine neue Pflanzenart erzeugt werden. Die beiden
Ausgangsspezies sind n = 4 und n = 8. Wie viele Chromosomen erwarten
sie in den somatischen Zellen der neuen Art?}
\begin{todolist}% %
\item 6 
\item 12 
\item 24  
\item 48
\item anderes
\end{todolist}
\subsection*{Was ist ein Diplont?}
\begin{todolist}% %
\item Ein Organismus, in dessen Lebenszyklus es keine haploiden Zellen
gibt.  
\item Ein Organismus, in dessen Lebenszyklus es keine diploiden Zellen
gibt.  
\item Ein Organismus, in dessen Lebenszyklus es keine Mitose gibt.  
\item Ein Organismus, in dessen Lebenszyklus es keine Meiose gibt.  
\item Ein Organismus, in dessen Lebenszyklus die Produkte der Meiose in
die Befruchtung eingehen.  
\end{todolist}

\subsection*{Wenn Sie Crossover bzw. den Austausch von genetischem Material in
einer Tetrade zwischen Chromatiden berücksichtigen, was wird in der
Meiose I voneinander getrennt?}
\begin{todolist} % %
\item Vollständige homologe Chromosomen.
\item Schwesterchromatiden von Zwei-Chromatid-Chromosomen.  
\item Chromosomen mit gleicher Genausstattung.  
\item Nur Telomere.  
\item Nur Centromere.  
\item Die Stränge einer doppelsträngigen DNA in zwei Einzelstränge.  
\end{todolist}
\subsection*{Eine Pflanze mit dem Genotyp AABbcc ist ...}
\begin{todolist} % %
\item homozygot in Bezug auf zwei Loci.
\item heterozygot in Bezug auf zwei Loci.  
\item rezessiv in Bezug auf drei Loci.  
\item Alle drei Angaben stimmen.  
\end{todolist}
\subsection*{Manchmal kann ein Gen die Auspägung eines anderen Gens verhindern
oder überdecken. Diese genetische Wechselwirkung bezeichnet man als}
\begin{todolist}% %
\item Genregulation.
\item Dominanz.
\item Epistasie.
\item zweite Paarregulation.
\end{todolist}

\begin{multicols}{2}
[\subsection*{Vater und Mutter haben den Genotyp AaBbCcDdEE. Die Loci liegen alle
auf verschiedenen Chromosomen. Welcher Anteil der Nachkommen hat
(was die Loci A, B, C, D, und E angeht) denselben Genotyp wie die
Eltern?}]
\begin{todolist}
\item $75 \% (3/4)$
\item $25 \% (1/4)$
\item ca. $14 \% (9/64)$
\item $6,25 \% (1/16)$
\item ca. $0,4 \% (1/256)$
\item 0
\end{todolist}
\columnbreak
Berechnung: 
\end{multicols}

\subsection*{Wenn alle Nachkommen einer Kreuzung heterozygot sind, sind die
Genotypen der Eltern wahrscheinlich ...}
\begin{todolist} % %
\item  AA, aa  
\item Aa, Aa  
\item Aa, aa 
\end{todolist}
\subsection*{Pleiotropie ist eine Genwirkung, bei der ...}
\begin{todolist}% %
\item viele Gene auf ein Merkmal wirken.  
\item mehrere Allele gleichzeitig exprimiert werden.  
\item beide Allele eines Locus die Ausprägung eines Merkmals
beeinflussen.  
\item ein Gen auf viele Merkmale wirkt.
\item ein Allel die Expression des anderen Allels vollständig maskiert.  
\end{todolist}
\subsection*{Welches der folgenden Paare kann KEIN Kind mit der Blutgruppe 0
haben? Angegeben sind die Blutgruppen der Eltern (Phänotypen).}
\begin{todolist} % %
\item  A x A
\item  B x B
\item  0 x AB
\item Keines der drei Paare.
\end{todolist}
\subsection*{Eine Fruchtfliege, die heterozygot für graue Körperfarbe und normale
Flügel ist (phänotypisch Wildtyp), wird mit einer schwarzen Fliege mit
Stummelflügeln gekreuzt. Die Verteilung der Phänotypen unter den
Nachkommen ist: 778 Wildtyp, 785 schwarz/stummelflügelig, 162 grau/
stummelflügelig, 158 schwarz/normalflügelig. Wie gross ist der Abstand
der beiden betrachteten Gene in [cM] auf der Basis dieser Kreuzung?}
\[
   \text{Gesamtzahl der Nachkommen} =
   \]
   \[
   \text{Rekombinanten Phänotypen} = 
   \]
   \[
   \text{Rekombinierungsfrequenz} = 
   \]
   \[
   \text{Abstand in cM} = 
   \]
\begin{todolist} % %
\item 0,17
\item 9
\item 17
\item 320
\item Zu wenig Informationen 
\end{todolist}
\subsection*{Thomas Bruder leidet an PKU, einer rezessiven Erbkrankheit. Die Eltern
des Bruders haben keine PKU. Wie wahrscheinlich ist es, das Thomas,
der in Bezug auf dieses Merkmal gesund ist, PKU übertragen kann?}
\[
\begin{array}{|c|c|}
\hline
\text{Elter 1} & \text{Elter 2} \\
\hline
\text{A} & \text{A} \\
\text{A} & \text{a} \\
\text{a} & \text{A} \\
\text{a} & \text{a} \\
\hline
\end{array}
\]
\[
\text{Genotypen: } \text{AA}, \text{Aa}, \text{aa}
\]
Die Wahrscheinlichkeiten sind:
\[
\text{P(AA)} = 
\]
\[
\text{P(Aa)} = 
\]
\[
\text{P(aa)} = 
\]
Da Thomas gesund ist, kann sein Genotyp entweder \( \text{AA} \) oder \( \text{Aa} \) sein. Die Wahrscheinlichkeit, dass Thomas ein Träger (\text{Aa}) ist, beträgt:
\[
\text{P(Träger)} = \text{P(Aa)} =  
\]
\begin{todolist} % %
\item $25 \% (1/4).$
\item $33,3 \% (1/3).$
\item  $50 \% (1/2).$
\item $66,6 \% (2/3).$
\end{todolist}
\subsection*{Die an einem RNA-Template (Matrize) synthetisierte DNA nennt man}
\begin{todolist}% %
\item Restriktionsenzym  
\item Reverse Transkriptase  
\item  cDNA
\item Palindrom  
\item Plasmid  
\end{todolist}
\subsection*{Eine Frau ist rotgrün-blind. Welchen 'Phänotyp' (bezogen auf rotgrün-
Blindheit) hat ihr Vater?}
\begin{todolist}% %
\item (Die Information reicht nicht aus, um das zu bestimmen.)
\item Er kann normal Farben sehen. 
\item  Er ist rotgrün-blind.  
\item Er hat zwei Y-Chromosomen.
\end{todolist}
\subsection*{Wenn das erste Kind eines Paares ein Junge war, wie hoch ist dann die
Wahrscheinlichkeit, dass das nächste Kind ein Mädchen sein wird?}
\begin{todolist}% %
\item $100\%$
\item $50\%$   
\item $33\%$
\item $25\%$
\item $0\%$
\end{todolist}
\subsection*{Bei einer Katzenvarietät ist das Auftreten von Silberflecken in der Iris der
Augen an das Geschlecht gekopppelt. FF-, Ff-, oder FY-Katzen haben
Flecken, ff- oder fY-Katzen haben keine. Daraus können wir schließen,
dass ....}
\begin{todolist}% %
\item sämtliche männliche Nachkommen von Weibchen, die Flecken
aufweisen, ebenfalls Flecken haben werden.  
\item ein Weibchen, das keine Flecken hat, nur männliche
Nachkommen ohne Flecken haben wird.
\item weibliche Katzen gefleckt oder nicht gefleckt sein können, Kater
dagegen nie Flecken haben.  
\end{todolist}
\subsection*{Bei vielen Vögeln haben die Männchen den Genotyp ZZ, und die
Weibchen den Genotyp ZW. Welches Elternteil bestimmt in diesem
System, welches Geschlecht die Nachkommen haben?}
\begin{todolist}% %
\item das Männchen
\item das Weibchen  
\item keines der beiden Elternteile
\item das Weibchen bestimmt das Geschlecht der weiblichen
Nachkommen, das Männchen das der männlichen
\end{todolist}

\subsection*{Bei der Erstellung von genetischen (Gen-) Karten werden aus Daten zu
Rekombinationsfrequenzen zwischen Genen genetische Distanzen
errechnet und die Abstände in cM angegeben. Was ist die Ursache für
dafür, dass Rekombinationsfrequenz-Werte (in $\%$) nicht einfach additiv in
Centi-Morgan umgerechnet werden können?}
\begin{todolist}% %
\item Weil Gene oft auf verschiedenen Kopplungsgruppen liegen.  
\item Weil Einzelcrossover-Ereignisse zu Rekombinanten führen, die
gegenüber den Parentaltypen einen Austausch zeigen.  
\item Weil Gene im Genom öfter die Position verändern.  
\item Weil Doppelcrossover-Ereignisse zwischen zwei Genen die
Rekombinationsfrequenz-Werte reduzieren.  
\end{todolist}

\subsection*{Welche Bedingungen muss eine Population bzw. die Allele in der
Population erfüllen, damit das Hardy-Weinberg-Gesetz für die
Allelfrequenzen in Populationen gültig ist?}
\begin{todolist}% %
\item Alle Allele eines Gens/Lokus müssen mit gleicher Häufigkeit auftreten.  
\item Manche Allele müssen wegen einer negativen Auswirkung auf die
Reproduktionsrate weniger häufig an die Nachfolge-Generation
weitergegeben werden.  
\item Die Population muss unbeschränkte, zufallsbedingte genetische
Durchmischung aufweisen.  
\end{todolist}

\subsection*{Was für eine Mutation (welche das/die fragliche(n) Gen(e) betrifft) wäre
wahrscheinlich am wenigsten schädlich für das zu bildende Protein?}
\begin{todolist}% %
\item Substitution einer Base durch eine andere   
\item Insertion einer Base  
\item Deletion einer Base  
\item Inversion eines Teilabschnitts des Chromosoms  
\item Alle genannten Mutation wären wahrscheinlich gleich schädlich  
\end{todolist}

\subsection*{Den Austausch von Chromosomenabschnitten zwischen nicht-homologen
Chromosomen bezeichnet man als ...}
\begin{todolist}% %
\item Inversion.  
\item Transduktion.  
\item Translokation. 
\item Transformation.  
\end{todolist}

\subsection*{Eine Person, die den Genotyp XXY aufweist, hat ....}
\begin{todolist}% %
\item eine Chromosomenaberration.  
\item eine autosomal rezessive Krankheit.   
\item eine autosomal dominante Krankheit.  
\item ein geschlechtsgekoppeltes rezessives Merkmal.  
\end{todolist} 

\subsection*{Ein Mensch mit Down-Syndrom hat ...}
\begin{todolist}% %
\item ein seltenes autosomal-rezessives Merkmal.  
\item ein zusätzliches Chromosom 21.
\item ein spezifisches Gen, das zu einer geistigen Behinderung führt.  
\item ein zusätzliches Chromosom 21 sowie ein spezifisches Gen,
das zu einer geistigen Behinderung führt.  
\end{todolist}

\subsection*{Welche der folgenden Konstellationen trifft für die einzige bekannte lebensfähige Monosomie beim Menschen zu?}
\begin{todolist}% %
\item XXY  
\item XYY  
\item XY  
\item Y0  
\item X0     
\end{todolist}

\subsection*{Wie viele (protein-kodierende) Gene befinden sich in etwa auf dem Y-Chromosom des Menschen?}
\begin{todolist}% %
\item Keines
\item ca. 70
\item etwa 700
\item Tausende
\end{todolist}

\subsection*{Die snRNPs genannten Partikel sind ...}
\begin{todolist}% %
\item Bestandteile des Spleißosoms. 
\item am Entfernen von Exons aus dem Primärtranskript beteiligt.  
\item ein bestimmter Typ spezialisierter Kohlenhydrate.  
\item ein Zwischenprodukt beim Zusammenbau des Ribosoms.  
\item an der Einschleusung von Proteinen ins ER beteiligt.  
\end{todolist}

\subsection*{Definieren Sie 'Spleißosom'}
\bigskip\bigskip\bigskip\bigskip\bigskip\bigskip
\subsection*{Normalerweise werden bei eukaryotischen Genen ...}
\begin{todolist}% %
\item die Exons nicht transkribiert.  
\item die Introns nicht transkribiert.  
\item die Introns zwar in RNA transkribiert, aber diese RNA verlässt den
Zellkern nicht.  
\item die Exons zwar in RNA transkribiert, aber diese RNA verlässt den
Zellkern nicht.  
\item Exons und Introns in RNA transkribiert, und diese RNA verlässt den
Zellkern.  
\end{todolist}

\subsection*{Beispiele für Totipotenz sind:}
\begin{todolist}% %
\item Mutationen, die zur Bildung eines Organs an Stelle eines anderen führen,  
\item die Regeneration einer kompletten Pflanze aus einzelnen Zellen des
Blattes oder der Wurzel,  
\item die Teilung embryonaler Zellen,  
\item die Abspaltung der Keimbahnzellen früh in der Embryogenese.  
\end{todolist}
\subsection*{Definieren Sie 'Totipotenz'}
\bigskip\bigskip\bigskip\bigskip\bigskip\bigskip


\subsection*{Eine einzelne embryonale Stammzelle kann sich nicht zu einem
kompletten Embryo entwickeln, weil ...}
\begin{todolist}% %
\item bestimmte Gene fehlen, die durch Rekombination verloren gegangen
sind.  
\item alle Meisterkontrollgene endgültig abgeschaltet sind.   
\item das Genom als Heterochromatin vorliegt.  
\item Informationen fehlen, die aus dem Cytoplasma der Eizelle stammen.    
\end{todolist}

\subsection*{Von einer Möhre werden durch Klonierung (Regenerierung ganzer
Pflanzen aus Zellen einer Pflanze) Nachkommen ("Klone") erzeugt.
Diese Nachkommen werden in einem grossen Feld angezogen und
....}
\begin{todolist}% % 
\item sehen alle absolut identisch aus, weil sie das gleiche Genom
enthalten.  
\item sehen alle absolut identisch aus, weil beim Klonieren Mutationen
ausgeschlossen sind.
\item sehen alle ähnlich aber im Detail doch unterschiedlich aus, weil sich
kleine Umwelteinflüsse auf das Erscheinungsbild von Pflanzen
auswirken.
\end{todolist}

\subsection*{Wodurch werden die Körperachsen in der frühen Entwicklung von D.
melanogaster festgelegt?}
\begin{todolist}% %
\item Durch Konzentrationsgradienten von Genprodukten.
\item Durch den Ort der Befruchtung der Eizelle durch die Samenzelle.
\item Durch Morphogene, die die Samenzelle in die Zygote einbringt.
\item Durch die Schwerkraft, die Stärkekörner zu einer Seite der Zygote
zieht.
\end{todolist}

\subsection*{Apoptose ist ...}
\begin{todolist}% %
\item eine Art biologischer Computervirus, der die Aktivität von Zellen
lahmlegt.  
\item ein Medikament, das zielgerichtet bestimmte Zellen abtötet.  
\item ein Mechanismus, der durch Aktivierung eines bestimmten
Programms von Proteasen und Nukleasen den Zelltod auslöst.  
\item ein Prozess, der normalerweise in Krebszellen abläuft und der die
Krebsentwicklung fördert.  
\end{todolist}

\subsection*{Bei Eukaryoten bezeichet man bestimmte Strukturen im Kern, die aus ca.
200 bp DNA und Protein bestehen, als ....}
\begin{todolist}% %
\item Histosomen  
\item Ribosomen  
\item Genetosomen  
\item Centriolen  
\item Nucleosomen
\end{todolist}
 
\subsection*{Was ist für den Unterschied in der Genomgrösse zwischen Arabidopsis
thaliana und Zea mays die wesentliche Ursache?}

\begin{todolist}% %
\item Monocotelydonen enthalten mehr Gene als Dikotelydonen  
\item Dikotelydonen enthalten mehr Gene als Monocotelydonen.  
\item Die Grösse der erwachsenen Pflanze.  
\item Repetitive DNA
\item Längere Introns.  
\end{todolist}
\subsection*{Definieren Sie 'Monocotelydonen und Dikotelydonen'}
\bigskip\bigskip\bigskip\bigskip\bigskip\bigskip


\subsection*{Wie kann man sich die Ausbreitung von Retrotransposons in Genomen
erklären?}
\begin{todolist}% %
\item Durch Fehler bei der Rekombination zwischen homologen
Chromosomen.  
\item Durch die Übernahme von retroviralen Funktionen wie reverse
Transkription und Insertion der entstandenen cDNA
\item Durch eine Abfolge von Inversionen und anderen
Chromosomenmutationen im Verlauf der Evolution  
\item Durch "cut and paste"-Mechanismen auf der Ebene der DNA wie bei
Transposons.  
\end{todolist}

\subsection*{Erklären Sie die Ausbreitung von Retrotransposons}
\bigskip\bigskip\bigskip\bigskip\bigskip\bigskip



\subsection*{Die Embryonalentwicklung von Pflanzen und Tieren ...}
\begin{todolist}% %
\item hat gemeinsam, das in beiden Fällen Zell- und Gewebebewegungen 
von Bedeutung sind.  
\item ist stark von Umweltreizen abhängig.  
\item hat gemeinsam, das es zentrale Kontrollgene gibt, die für DNA-
bindende Proteine kodieren.  
\item Keine der vorherigen Aussagen ist richtig.  
\end{todolist}
\subsection*{Wie schreibt man DNA-Sequenzen?}


\begin{center}
    \textit{(ss-DNA)}
\end{center}
  bzw. \newline

 \begin{center}
    \textit{(ds-DNA)}
\end{center}


\subsection*{Wohin wandert DNA im elektrischen Feld?}

\bigskip\bigskip\bigskip


\subsection*{Welche strukturellen Eigenschaften haben DNA und Polyphosphat gemeinsam? }
\begin{itemize}
    \item 
    \item 
    \item 
    \item 
\end{itemize}

\subsection*{Wie färbt man DNA?}
\begin{itemize}
    \item 
    \item 
\end{itemize}

\subsection*{Warum ist G-C stabiler als A-T?}
\bigskip\bigskip\bigskip

\begin{itemize}
    \item Guanin (G) und Cytosin (C) ...
    \item Adenin (A) und Thymin (T) ...
\end{itemize}

\subsection*{Wie unterscheidet sich DNA von RNA?}
\begin{itemize}
    \item 
    \begin{itemize}
        \item DNA  
        \item RNA  
    \end{itemize}
    \item 
    \begin{itemize}
        \item DNA  
        \item RNA  
    \end{itemize}
    \item 
    \begin{itemize}
        \item DNA  
        \item RNA 
    \end{itemize}
    \item 
    \begin{itemize}
        \item DNA: 
        \item RNA: 
    \end{itemize}
    \item 
    \begin{itemize}
        \item DNA 
        \item RNA 
    \end{itemize}
    \item 
    \begin{itemize}
        \item DNA: 
        \item RNA: 
    \end{itemize}
\end{itemize}

\subsection*{Warum ist DNA stabiler als RNA?}
\bigskip\bigskip\bigskip\bigskip

\begin{itemize}
    \item 
    \item 
    \item 
    \item 
\end{itemize}

\subsection*{Aufgabe: Die kodierende Region eines Gens weist bezüglich der Basenzusammensetzung 5\% Adenin und 45\% Guanin auf}
Erinnerung: \newline
\centerline{$A+T+G+C=100\%$}
\centerline{$A = T$}
\centerline{$C = G$}
\subsubsection*{1. Was ist der prozentuale Anteil an Cytosin?}
\bigskip\bigskip
\subsubsection*{2. Was ist in der Entsprechenden RNA der prozentuale Anteil an Uracil?} 
\bigskip\bigskip
\subsubsection*{3. Was ist der prozentuale Anteil an Purinen? } 
\bigskip\bigskip\bigskip\bigskip\bigskip\bigskip


\subsection*{Welcher Phosphatrest wird in die DNA eingebaut?}
\bigskip\bigskip\bigskip\bigskip


\subsection*{Welcher Phosphatrest wird bei Kinase-Reaktionen transferiert?}
\bigskip\bigskip\bigskip\bigskip

 
\subsection*{Was kann man über die Kinetik der DNA-Synthese sagen?}
\begin{itemize}
    \item 
    \item 
    \item 
    \item 
    \item 
    \item
\end{itemize} 

\subsection*{DNA-Replikation}
\subsubsection*{1. Ist wie gerichtet?} 
\bigskip\bigskip\bigskip\bigskip\bigskip\bigskip

\subsubsection*{2. Verhält sich wie bezüglich der Ausgangs-DNA und Tochter-DNA?} 
\bigskip\bigskip\bigskip\bigskip\bigskip\bigskip

\subsubsection*{3. Verhält sich wie bezüglich des Ablaufs der Replikation an den beiden DNA-Strängen?} 
Es gibt einen führenden und einen verzögerten Strang.
\begin{table}[H]
    \centering
    \begin{tabular}{|p{7cm}|p{7cm}|} \hline
     führend & verzögert \\ \hline \hline
    \bigskip\bigskip\bigskip &   \bigskip\bigskip\bigskip \\\hline
\end{tabular}
\end{table} 



\subsection*{Geben Sie die Werte der Parameter der DNA-
Superhelikalitat an? Wie erfolgen die Umwandlungen?}
\subsubsection*{Formel von LK, T und W:}
\bigskip\bigskip
\begin{figure}[H] 
    \centering
    \includegraphics[width=10cm]{bilder/linkerausfüllen.png}
    \label{fig:my_label}
    \caption{}
\end{figure}
\subsection*{Aufgabe: In einer Probe seien DNA und RNA enthalten. Sie sollen nur die RNA zerstören. Wie gehen Sie vor?}
\bigskip\bigskip\bigskip\bigskip


\subsection*{Welche(n) Parameter der Superhelikalität der DNA können nur Topoisomerasen verändern?}
\begin{todolist}
\item  Linking Number (LK)
\item Verdrilling Twist (T) 
\item Anzahl der Überkreuzungen Writhe (W)
\end{todolist}
\subsection*{Wie unterscheidet man superhelikale von 'entspannter' DNA?}
\begin{table}[H]
    \centering
    \begin{tabular}{|c||p{5cm}|p{5cm}|} \hline 
         & Superhelikal & Entspannt  \\ \hline\hline
        Verdrillung&& \\ \hline
        Topologische Eigenschaften&& \\ \hline
        Funktionelle Bedeutung && \\ \hline
    \end{tabular}
    \label{tab:my_label}
\end{table} 
\subsection*{Lückentext}
Bei der Replikation der DNA von E. coli arbeitet die DNA-Polymerase III. sehr genau. Sie synthetisiert nur jedes \textbf{.........................}te Basenpaar falsch, weil sie selbst über eine Reparaturfunktion verfügt, die \textbf{.........................}  ist. \newline
Replizierte DNA weist jedoch noch weniger Fehler auf: nur jedes \textbf{.........................}  te Basenpaar falsch, weil es zusätzliche Reparaturmechanismen gibt, z.B. die \textbf{.........................} .

\subsection*{Wie können Mutationen induziert werden?}
\subsubsection*{1. Physikalisch, z.B.} 
    
    \begin{itemize}
    \item 
    \item 
    \item
\end{itemize}
\subsubsection*{2. Chemisch, z.B.} 
\bigskip\bigskip
\begin{itemize}
    \item 
    \item 
    \item 
\end{itemize}

\subsection*{Aufgabe: Bei der Entschlüsselung des genetischen Codes wurden RNA-Moleküle  die entweder nur aus U oder nur aus CU zusammengesetzt waren, translatiert und die entstehenden Proteine charakterisiert. Welche Aminosäure(n) werden durch eine RNA kodiert, die sich durch aufeinander folgende AG's zusammensetzt?}
Was ist hier gefordert?
\begin{enumerate}
    \item Codons bilden  \newline \bigskip\bigskip
    \item Codons zu Aminosäuren zuordnen \newline \bigskip\bigskip
\end{enumerate} 
\subsection*{Aufgabe: Ein Protein hätte drei Aminosäuren (1.Methionin, 2. Tryptophan, 3. Tryptophan). Schreiben Sie die Protein- und die entsprechende mRNA-Sequenz auf! Benennen Sie
die vorderen und hinteren Enden! Benennen Sie den Vorgang}
\begin{figure}[H]
    \centering
    \includegraphics[width=10cm]{bilder/sonne.png}
    \caption{}
    \label{fig:my_label}
\end{figure}
Was ist hier gefordert? 
\begin{enumerate}
    \item Bestimmung der Proteinsequenz \newline \bigskip\bigskip
    \item Bestimmung der Codons durch Ablesen von der Sonne \newline \bigskip\bigskip\bigskip\bigskip 
    \item Bestimmung der RNA-Sequenz aus den Codons \newline \bigskip\bigskip\bigskip\bigskip
    \item Benennen wo Ende und Anfang sind \newline \bigskip\bigskip\bigskip\bigskip
    \item Benennung des Vorgangs \newline \bigskip\bigskip
\end{enumerate}



\subsection*{Für ein E. coli-Gen sei mit einer molekular biologischen Technik der Transkriptionsstartpunkt (das erste Nukleotid der mRNA) bestimmt worden und in der angegebenen Sequenz fett markiert. Identifizieren Sie die Promotor-Elemente!}
\centerline{AAAAAATCTTGACACCCAGCATCAGTAGACATAATATGACCACA\textbf{T}GGACCAAA} 
\noindent Was wird hier gefordert? 
\subsubsection*{Konsensussequenzen der E. coli Promotoren finden}
\bigskip\bigskip\bigskip

\subsection*{Die RNA-Polymerase ist wesentlich für die Transkription in E. coli. Welchem Enzym der DNA-Replikation (Primase oder DNA-Polymerase III) ähnelt sie mehr?}
Primase und RNA-Polymerase. 
\begin{table}[h] 
    \centering
    \begin{tabular}{|p{3cm}||p{4cm}|p{4cm}|p{4cm}|}\hline
         & Primase &  DNA-Polymerase III & RNA-Polymerase\\ \hline \hline
         Aktivität /was wird synthetisiert & & &  RNA \\\hline
         Vorläufer/Substrate &  & & NTPs \textit{(Ribonukleosidtriphosphate)} \\\hline
         Matrize & & &  ssDNA \\\hline
         Starter oder Primer & & &  Kein Primer benötigt \\\hline
    \end{tabular}
    \label{tab:my_label}
\end{table}

\subsection*{Wie viele verschiedene Aminosäuren gibt es in
Proteinen?}
\bigskip\bigskip

\subsection*{Anzahl der während der Proteinbiosynthese (ko-translational) eingebauten Aminosäuren}
\bigskip\bigskip

\subsection*{Woraus setzt sich ein Ribosom in Bakterien zusammen?}
\begin{todolist}
\item nur aus RNA 
\item nur aus Protein
\item  aus RNA und Protein
\item  aus 3 RNA Molekülen und aus 1 Protein
\item aus 3 RNA Molekülen und aus vielen Proteinen 
\item  aus 2 Untereinheiten
\end{todolist} 

\subsection*{Ordnen Sie die
Translationsfaktoren EF-G,EF-TU, IF1, IF2, IF3, RF und RRF zeitlich in den Prozess der Translation ein und weisen Ihnen eine Funktion dabei zu. 
}
\begin{enumerate}
    \item  ........................................
    \begin{enumerate}
        \item ........................................
        \item ........................................
        \item ........................................
    \end{enumerate}
    \item ........................................
    \begin{enumerate}
        \item ........................................
        \item ........................................
    \end{enumerate}
    \item ........................................
    \begin{enumerate}
        \item ........................................
        \item ........................................
    \end{enumerate}
\end{enumerate}
\subsection*{Was ist  die Shine-Dalgarno-Sequenz?
Wo liegt sie? Was ist  ihre Funktion?}
\bigskip\bigskip \bigskip\bigskip\bigskip\bigskip

\subsection*{Wie sind Ribosomen aus Archaea und Eukaryoten zusammengesetzt?} 
\begin{table}[H]
    \centering
    \begin{tabular}{ |p{2cm}|| p{4cm}| p{4cm}| p{4cm}|}
    \hline
     & Archaea & Bacteria & Eukaryoten \\ \hline \hline
    Typ &   &  & \\ \hline
    zusammen gesetzt aus..  &   &  & \\ \hline
    rRNA &  &   &  & \\ \hline
    ribosomale Proteine  &   &  & \\ \hline
    Hemmstoffe &   &  & \\ \hline
\end{tabular}
\end{table}

\subsection*{Was ist das katalytisch aktive Prinzip des Ribosoms?}
\bigskip\bigskip\bigskip\bigskip

\subsection*{Aufgabe: Ordnen Sie die Vorgänge bei der homologen Rekombination. (nicht alle existieren)}
Branch migration Wandern der Holliday-Struktur Eindringen eines Einzelstrangs
Eindringen eines Doppelstrangs Auf losen der Holliday-Struktur Doppelstrang-Bruch
Ausbildung von 3'-Überhängen Ausbildung von 5'-Überhängen
\begin{enumerate}
    \item ........................................
    \item ........................................
    \item........................................
    \item ........................................
    \item ........................................
\end{enumerate}
Es existieren nicht: 
\begin{itemize}
    \item ........................................
    \item ........................................
    \item ........................................
\end{itemize}

\subsection*{Bei der Rekombination zweier homologer DNA-Regionen mit den nach einander liegenden Genen 1, 2, 3 und 4 bildet sich eine Holliday-Struktur aus, die über die Gene 2 und 3 hinweg wandert. Zeichnen Sie die Schritte der Rekombination und lösen Sie die Holliday-Struktur auf zwei Arten auf.} 
\bigskip\bigskip\bigskip\bigskip\bigskip\bigskip\bigskip\bigskip\bigskip\bigskip\bigskip\bigskip\bigskip\bigskip\bigskip\bigskip\bigskip\bigskip\bigskip\bigskip\bigskip\bigskip\bigskip\bigskip

\subsection*{Bei welchen Genen kann Genkonversion auftreten?}
\bigskip\bigskip\bigskip\bigskip

\subsection*{Als Oswald Avery zu einem Transformationsansatz mit einem Gemisch aus durch Hitze getöteten 5-Zellen und lebenden R-Zellen von Streptococcus pneumoniae folgende Substanz zusetzte unterblieb die Transformation und die mit dem Gemisch infizierte Maus überlebte.
}
\begin{todolist}
\item Lipase
\item Lipase + Protease 
\item Phosphatase + Protease
\item Phosphatase 
\item  DNase
\item RNase
\item Protease
\end{todolist}
\subsection*{Bei Bakterien gibt es verschiedene Möglichkeiten der DNA-Übertragung. Der direkte Zell-Zell-Kontakt ist wesentlich bei der ...}
\begin{todolist}
\item Transformation
\item  Konjugation
\item Transduktion
\item Phagen-Lyse 
\item Transkription
\item Lysogenie
\item Translation
\item DNA-Replikation
\end{todolist}
\subsection*{Sie transformieren einen Tryptophan-, Arginin- und Leucin- auxotrophen Bakterien-Stamm trp- arg- leu- mit DNA eines Arginin- und Leucin-prototrophen, aber Tryptophan-auxotrophen Stammes trp - arg + leu+. 
Welche Rekombinante(n) kann/können entstehen?}
\bigskip\bigskip\bigskip\bigskip

\subsection*{Enzyme und Proteine der DNA-Replikation}
\begin{table}[H]
    \centering 
    \begin{tabular}{|p{4cm}|p{11cm}|}
    \hline
     &  Initiation der Replikation am oriC\\\hline 
     & Bindet an einzelsträngige DNA  \\\hline
    Helikase &  \\\hline
    Primase &   \\\hline
     & Synthetisiert den neuen DNA-Strang \\\hline
     & Entfernt den RNA-Primer (5'-3'-Exonuklease) und füllt den Strang auf \\\hline
    Ligase &  \\\hline
    Topoisomerase &  \\\hline
\end{tabular}
\end{table}
\subsection*{tRNA}
\subsubsection*{Wie lang sind tRNA-Moleküle?}
\bigskip\bigskip
\subsubsection*{Welche Struktur nehmen tRNA-Moleküle ein?}
\bigskip\bigskip\bigskip\bigskip
\subsubsection*{Was kann man über die Basen in tRNA sagen?}
\bigskip\bigskip\bigskip\bigskip
\bigskip\bigskip\bigskip\bigskip

\subsubsection*{Welche Funktion hat tRNA?}
\bigskip\bigskip\bigskip\bigskip \textit{(Translation)}. 
\subsection*{Beschreiben Sie das trp-Operon und seine Kontrolle}
\begin{table}[H]
    \centering
    \begin{tabular}{|c||p{5cm}|p{5cm}|} \hline
         & Repression und Aktivierung
der Transkriptionsinitiation &  Transkriptionelle
Attenuation \\ \hline \hline
        Protein & Repressor oder Aktivator &   \\ \hline 
         Transkription & Initiation wird reguliert &   \\ \hline 
        Vorkommen & Prokaryoten und Eukaryoten &  \\ \hline 
    \end{tabular}
    \label{tab:my_label}
\end{table}
\subsection*{Sie wollen ein Protein mit einem Wirtsbakterium produzieren. Dazu klonieren
Sie das Gen für dieses Protein mitsamt der für Transkription, Translation und Sekretion notwendigen Sequenzen. Welches Bakterium wählen Sie?}
\bigskip\bigskip\bigskip\bigskip


\subsection*{Bei der Klonierung von Genen nutzt man gerne die a-Komplementation.
Wieviele Gene dafür sind wo kodiert?}
\bigskip\bigskip\bigskip\bigskip\bigskip\bigskip\bigskip\bigskip\bigskip\bigskip\bigskip\bigskip

\subsection*{Wo und warum setzt man die Taq Polymerase ein?}
\bigskip\bigskip\bigskip\bigskip
\subsection*{Beschreiben Sie Taq Polymerase bezüglich ...}
\subsubsection*{... der Substrate}
\bigskip\bigskip
\subsubsection*{... der Matrize}
\bigskip\bigskip
\subsubsection*{... der Aktivität}
\bigskip\bigskip
\subsection*{Lückentext}
Bei der Expressionsklonierung können Sie experimentell der Frage nach dem ................... wissenschaftsphilosophischen Kriterium nachgehen. 
\subsection*{Welche Elemente muß die zu klonierende DNA aufweisen?}
\begin{itemize}
    \item 
    \item 
    \item 
    \item 
\end{itemize}
\subsection*{Beschreiben Sie die Schritte der Gram-Färbung!}
Beschreiben Sie die Schritte der Gram-Färbung!
\begin{enumerate}
    \item 
    \item 
    \item 
    \item 
\end{enumerate}

\subsection*{SSU rDNA-Sequenzanalysen erlauben es alle Lebewesen der Welt phylogenetisch zu betrachten. Was ist die Einheit als evolutionärer Abstandsmesser, was als evolutionäre Uhr?}
\textbf{“Evolutionäre Abstandsmesser”} \newline \bigskip\bigskip\bigskip\bigskip \newline
\textbf{“Evolutionäre Uhr”} \newline \bigskip\bigskip\bigskip\bigskip
\subsection*{Der phylogenetische Stammbaum der Lebewesen basiert auf SSU rDNA-Sequenzanalysen. Beschriften sie die 3 Reiche und zeichnen Sie die ungefähren Positionen von E. coli, Mensch, Maus und Mais ein!}
\begin{figure}[H]
    \centering
    \includegraphics[width=10cm]{bilder/maisposition.png}
    \caption{}
    \label{fig:my_label}
\end{figure}
\subsection*{Was katalysiert Lysozym?}
\bigskip\bigskip\bigskip
\subsection*{Was katalysiert Transpeptidase?}
\bigskip\bigskip\bigskip

\subsection*{Was katalysiert Transglykosidase?} 
\bigskip\bigskip\bigskip

\subsection*{Wie wirkt Penicillin (und andere beta-Lactam-Antikiotika)?}
\begin{itemize}
    \item 
    \item 
    \item 
\end{itemize}
\subsection*{Skizzieren Sie die Zellwand ...}
\subsubsection*{... eines Gram-positiven Bakterium}
\bigskip\bigskip\bigskip\bigskip\bigskip\bigskip
\subsubsection*{... eines Gram-negativen Bakteriums}
\bigskip\bigskip\bigskip\bigskip\bigskip\bigskip

\subsection*{Was versteht man
in der mikrobiellen Ökologie unter „great plate
count anomaly“?}
\bigskip\bigskip\bigskip\bigskip\bigskip\bigskip
\subsection*{Wie setzt sich das bakterielle Murein
zusammen? Wie ist es verknüpft?}
\subsubsection{Zusammensetzung}
\begin{itemize}
    \item 
    \item 
\end{itemize}
\subsubsection{Verknüpfung}
\begin{itemize}
    \item 
    \item 
    \item
\end{itemize}
\subsection*{Resistenz gegen Penicillin}
\subsubsection*{Wie kann ich Penicillin-resistente Bakterien
selektionieren?}
\bigskip\bigskip\bigskip
\subsubsection*{Welches Enzym kann Penicillin-Resistenz
vermitteln?}
\bigskip\bigskip\bigskip
\subsection*{Gel-Elektrophorese}
\subsubsection*{Welchen Typ der Gel-Elektrophorese nutzt
man zur Analyse und Größenfraktionierung
von DNA?}
\bigskip\bigskip\bigskip
\subsubsection*{Welchen Typ der Gel-Elektrophorese nutzt
man zur Analyse und Größenfraktionierung
von Proteinen?}
\bigskip\bigskip\bigskip\bigskip\bigskip\bigskip
\subsection*{Was ist das Prinzip der DNA-
Sequenzierung nach Sanger?}
\bigskip\bigskip\bigskip\bigskip\bigskip\bigskip
\subsection*{Beschreiben Sie die
Genaktivierung durch NtrC~P}
\begin{enumerate}
    \item 
\item 
\item 
\end{enumerate}

\subsection*{Listen Sie die Aminosäuren}
\bigskip\bigskip\bigskip
\bigskip\bigskip\bigskip
\bigskip\bigskip\bigskip




\end{document}